\section{Stationary endpoint roofs at a positive critical entrance}
\label{sec:v104-stationary}
The preceding comparison can be verified on a class for which the
endpoint speed has no positive lower bound. The condition on the
outgoing trace is an integral condition. Neither a derivative of that
trace nor a pointwise lower bound for the positive minimum is used.

Consider an oriented angular arm $x=\sigma(\theta-\theta_*)>0$,
$\sigma\in\{-1,1\}$, on one actual pairing sheet
with a regular nongrazing critical entrance. The incoming roof is
$a_*$ and, by \eqref{eq:v104-critical-input}, its weighted trace is
\begin{equation}\label{eq:v104-incoming}
 f_-^\varepsilon(x)=\frac{N_*}{Q_*(x)},\qquad
 \lambda_\circ\le Q_*(x)\le\Lambda_\circ,
 \qquad N_*\ge0.
\end{equation}
Suppose the outgoing roof on this arm is
\begin{equation}\label{eq:v104-profile}
 b(x)=a_*+d+\psi(x),\qquad 0\le d\le\delta/2,
\end{equation}
where $\psi$ is continuous, strictly increasing, positive for $x>0$,
and tends to zero at the left endpoint. Assume that its levels
$\delta/2,2\delta,3\delta$ are reached at points
$x_1<x_2<x_3$ inside this same physical arm. Set
\[
 I=(0,x_1),\qquad J=(x_2,x_3).
\]
The existence of $x_3$ is an actual angular-reach condition. It is
not inferred from an algebraic continuation beyond a true endpoint.

If $N_*>0$, put
\begin{equation}\label{eq:v104-fraction}
 w(x)=\min\{1,f_+^\varepsilon(x)/f_-^\varepsilon(x)\},
 \qquad \mu_J=\frac1{|J|}\int_Jw(x)\,dx.
\end{equation}
Thus $c^\varepsilon=f_-^\varepsilon w$ exactly. The fraction may
vanish on subsets of positive angular measure. If $N_*=0$, the
whole paired source on the arm is zero and is treated separately.

\begin{lemma}[Width and positive-trace balance]
\label{lem:v104-arm-balance}
For $N_*>0$ and $\mu_J>0$,
\begin{equation}\label{eq:v104-geometric-balance}
 \frac{\int_Ic^\varepsilon dx}{\int_Jc^\varepsilon dx}
 \le\frac{\Lambda_\circ}{\lambda_\circ\mu_J}\frac{x_1}{x_3-x_2}.
\end{equation}
If, in addition, $\psi$ is locally absolutely continuous and
\begin{equation}\label{eq:v104-elasticity}
 \frac{x\psi'(x)}{\psi(x)}\le K
           \quad\text{for almost every }x\in(x_2,x_3),
           \qquad K>0,
\end{equation}
then
\begin{equation}\label{eq:v104-arm-balance}
 \frac{\int_Ic^\varepsilon dx}{\int_Jc^\varepsilon dx}
 \le\frac{\Lambda_\circ K}
          {\lambda_\circ\mu_J\log(3/2)}.
\end{equation}
No lower bound on $\psi'$ or bound on $\partial_x\log\psi'$ is
part of either assertion.
\end{lemma}
\begin{proof}
On $I$ one has $c^\varepsilon\le N_*/\lambda_\circ$; on $J$,
$c^\varepsilon\ge N_*w/\Lambda_\circ$. Integration proves
\eqref{eq:v104-geometric-balance}. Since $\psi\ge2\delta$ on
$[x_2,x_3]$, its logarithm is absolutely continuous there. Hence
\[
 \log(3/2)=\int_{x_2}^{x_3}\frac{\psi'}{\psi}\,dx
                \le K\log(x_3/x_2).
\]
It follows that
\[
 x_3-x_2\ge x_2\bigl((3/2)^{1/K}-1\bigr)
              \ge x_1\log(3/2)/K.
\]
Substitution gives \eqref{eq:v104-arm-balance}. The numerator $N_*$
cancels; no positive lower bound for its size is required.
\end{proof}

\begin{theorem}[A stationary-arm essential-height estimate]
\label{thm:v104-stationary}
Take any disjoint family of actual angular arms with
\eqref{eq:v104-incoming}--\eqref{eq:v104-profile}. Suppose their
comparison intervals satisfy either
\[
 \frac{\Lambda_\circ}{\lambda_\circ\mu_J}
                    \frac{x_1}{x_3-x_2}\le\kappa,
\]
or the sufficient uniform bounds $K\le K_0$ and $\mu_J\ge\mu_0>0$
in Lemma~\ref{lem:v104-arm-balance}. Let $T_{H,\mathrm{arm}}^
\varepsilon$ be any original selected paired source on the union
of their source intervals $I$. At finite collision count,
\begin{equation}\label{eq:v104-arm-finite}
 T_{H,\mathrm{arm}}^\varepsilon(t)
 \le\frac{\kappa}{\delta}\int_t^{t+\delta}
                        \mathsf T_{4\delta}^\varepsilon(u)\,du
 \quad\text{for almost every }t.
\end{equation}
Consequently
\begin{equation}\label{eq:v104-arm-height}
 \mathcal H_M(T_{H,\mathrm{arm}}^\varepsilon)
 \le C\kappa\{\varepsilon^{1/16}
                                    +(1+\varepsilon^{-1})\delta\},
 \qquad
 \kappa=\frac{\Lambda_\circ K_0}
              {\lambda_\circ\mu_0\log(3/2)}
\end{equation}
in the sufficient case. Every original relative position $0<s<1$
is included. The family may have arbitrarily many physical words.
\end{theorem}
\begin{proof}
For a source angle $x\in I$ the residence is at most
$d+\delta/2\le\delta$. For a comparison angle $y\in J$, it is
less than $d+3\delta\le7\delta/2<4\delta$. Their incoming roofs
are equal, and
\[
 b(y)-b(x)>2\delta-\delta/2=3\delta/2>\delta.
\]
Thus the positive inclusion \eqref{eq:v104-inclusion} holds with
$r=\delta$. Lemma~\ref{lem:v104-arm-balance} proves the mass balance.
Different arms are disjoint in the original physical angular space;
therefore their comparison intervals have overlap one, even if
geometrical boundary points coincide on different incident sheets.
The proof of Theorem~\ref{thm:v104-block}, with every sign positive,
gives \eqref{eq:v104-arm-finite}. Its right side is at most
$2\kappa A_\delta\mathsf T_{4\delta}^\varepsilon$. The inherited
paired-average estimate proves \eqref{eq:v104-arm-height}.
Arms with $N_*=0$ have no positive minimum and add zero.
\end{proof}

The arm assumption, including the weighted fraction in
\eqref{eq:v104-fraction}, is imposed on actual endpoint traces.
It is weaker than a positive pointwise lower bound or logarithmic
variation bound on their minimum. Failure of this average condition
is quantified by the positive deficit in \eqref{eq:v104-deficit},
not assigned zero mass.

\subsection{Paying a deficient comparison trace by the original volume}
The average fraction in \eqref{eq:v104-fraction} need not be bounded
below when the original weighted trace variation is available. Its
possible loss has a direct pointwise payment. Retain the actual
uncut source $v^0$ from \eqref{eq:v104-uncut-input}.

\begin{theorem}[Stationary-arm height without a trace-fraction hypothesis]
\label{thm:v104-paid-arms}
On a disjoint family of the actual critical arms above, assume only
\eqref{eq:v104-incoming}--\eqref{eq:v104-profile} and a uniform
geometric width ratio
\[
 \frac{\Lambda_\circ}{\lambda_\circ}
              \frac{x_1}{x_3-x_2}\le\kappa.
\]
There is no lower bound on $\mu_J$. For every original selector
$0\le H\le1$, at finite count and almost every target roof,
\begin{align}
 T_{H,\mathrm{arm}}^\varepsilon(t)
 \le{}&\frac{\kappa}{\delta}\int_t^{t+\delta}
                       \mathsf T_{4\delta}^\varepsilon(u)\,du\notag\\
 &+C\kappa(1+\varepsilon^{-1})
                         \int_{t-\delta}^{t+4\delta}v^0(u)\,du.
                                           \label{eq:v104-paid-arms-finite}
\end{align}
The positive comparison deficit itself obeys
\begin{equation}\label{eq:v104-paid-deficit}
 \sum_j(1-\alpha_j)T_{H,j}^\varepsilon(t)
 \le C\kappa(1+\varepsilon^{-1})
                         \int_{t-\delta}^{t+4\delta}v^0(u)\,du.
\end{equation}
Consequently the entire selected arm source satisfies
\begin{equation}\label{eq:v104-paid-arms-height}
 \mathcal H_M(T_{H,\mathrm{arm}}^\varepsilon)
 \le C\kappa\{\varepsilon^{1/16}+(1+\varepsilon^{-1})\delta\}.
\end{equation}
The sufficient profile condition \eqref{eq:v104-elasticity} gives
$\kappa=\Lambda_\circ K_0/(\lambda_\circ\log(3/2))$ for $K\le K_0$.
No source endpoint-speed, incidence, residence or relative-position
lower cutoff is imposed.
\end{theorem}
\begin{proof}
For one arm set
\[
 P_j=\int_{J_j}f_-^\varepsilon(x)\,dx,\qquad
 E_j=\int_{J_j}[f_-^\varepsilon(x)-f_+^\varepsilon(x)]_+\,dx.
\]
The positive minimum gives the exact identity $D_j=P_j-E_j$.
The critical incoming formula and the geometric width ratio imply
$C_j\le\kappa P_j$. Therefore
\begin{equation}\label{eq:v104-deficit-variation}
                  (C_j-\kappa D_j)_+\le\kappa E_j.
\end{equation}
This is valid even when $D_j=0$. Arms with $N_*=0$ again contribute
zero; no division by a vanishing incoming trace is used.

The retained derivative bound for the original weighted density,
integrated on the same comparison segment, yields
\[
 [f_-^\varepsilon(x)-f_+^\varepsilon(x)]_+
 \le C(1+\varepsilon^{-1})
        \int_{a_j}^{b(x)} A(u,x)I_m(u,x)/z(u,x)\,du.
\]
One may first integrate on compact interior subintervals and pass to
the one-sided traces. The uncut volume input and Tonelli justify
the positive limiting integral. This pays a possible decrease of
the trace; it does not assume that the two traces are comparable.

At a roof with $T_{H,j}^\varepsilon(t)>0$, some original source
segment on $I_j$ contains $t$. Since every such segment has length
at most $\delta$ and has the same incoming roof $a_j$, one has
$t-\delta<a_j<t$. Each comparison segment on $J_j$ has length
less than $4\delta$. Its entire interior consequently lies in
$(t-\delta,t+4\delta)$. The comparison angular sets are disjoint
on the original physical sheet space. After their interiors are
summed, the Borel segment identity bounds their combined variation
by the integral of the same complete $v^0$ over this one roof window,
with multiplicity one. Equation~\eqref{eq:v104-deficit-variation}
and \eqref{eq:v104-deficit} prove \eqref{eq:v104-paid-deficit}.
Adding the paid comparison in Theorem~\ref{thm:v104-block} gives
\eqref{eq:v104-paid-arms-finite}. The first term uses only positive
roof translates, since the incoming roof is constant on each arm.

Finally take the collision limsup with $\delta,\varepsilon,\kappa$
fixed. The paired-average input estimates the first integral. The
uncut volume input estimates the second by
$C\kappa(1+\varepsilon^{-1})\delta$. Auxiliary spectral bands are
enlarged only afterwards. This proves \eqref{eq:v104-paid-arms-height}.
\end{proof}

Thus poor outgoing trace balance on these arms is not assigned to an
uncontrolled positive remainder: its entire shortage is estimated at
the original target roof by \eqref{eq:v104-paid-deficit}. What still
requires physical information is the level reach and width geometry,
and coverage of the complete remaining paired family by such arms or
other comparisons.

\begin{corollary}[Finite-order stationary profiles]
\label{cor:v104-finite-order}
Suppose on an actual arm
\[
 \psi(x)=\beta x^k(1+O(x^\eta)),\qquad
 \psi'(x)=k\beta x^{k-1}(1+O(x^\eta)),
 \qquad \beta,\eta>0,\quad k\ge1.
\]
For all sufficiently small widths on that arm,
\eqref{eq:v104-elasticity} holds with $K=k+1$. Provided the required
physical reach holds, Theorem~\ref{thm:v104-paid-arms}
applies also for $k>1$, when the endpoint speed vanishes at the arm's
initial point. No uniform bound on physical orders $k$, on the size
of the asymptotic neighborhood or on comparison reach is inferred
from finite-count definability.
\end{corollary}
\begin{proof}
The displayed asymptotics give $x\psi'/\psi\to k$. The assertion
then follows from Lemma~\ref{lem:v104-arm-balance} and
Theorem~\ref{thm:v104-paid-arms}.
\end{proof}

\subsection{Recovery can occur when every preceding clock is small}
This elementary pairing-sheet calculation distinguishes the new
comparison from admission by any of the three previous flows. It is
not asserted to be another physical Lorentz realization. Set
$a=0$, $b(x)=x^k$, $c=1$, $k\ge2$, and restrict the comparison
collar to $0<\ell<4\delta$. The incoming-distance clock is zero.
For the outgoing-distance flow at positive roof time $h$,
\[
 x_h=(x^k+h)^{1/k},\qquad
 J_h^+=\left(\frac{x}{x_h}\right)^{k-1}.
\]
Its forward contribution for $0<h<\delta$ is at most
\begin{equation}\label{eq:v104-model-clock}
 \frac{kx^{k-1}}{\delta}
        \bigl((x^k+\delta)^{1/k}-x\bigr).
\end{equation}
The negative-time contribution is at most $x^k/\delta$, because
the starting gap must remain smaller than the new residence. At
fixed relative position $s$, the relative clock has total roof reach
at most $4s\delta$, so its capped value is at most $4s$.

For any fixed $q>0$, take $s<q/8$ and then take
$x/\delta^{1/k}$ sufficiently small. Both terms in
\eqref{eq:v104-model-clock} and its negative-time bound tend to zero;
all three old clocks are below $q$ on a set of positive paired
measure. Nevertheless the intervals
\[
 I=(0,(\delta/2)^{1/k}),\qquad
 J=((2\delta)^{1/k},(3\delta)^{1/k})
\]
satisfy the positive comparison, with
$C_I/D_J\le k/\log(3/2)$. The new inequality controls their entire
source, including that positive-measure clock-failure set. Its proof
uses neither a velocity bound on $I$ nor passage of a flow through
$x=0$.
